% Options for packages loaded elsewhere
\PassOptionsToPackage{unicode}{hyperref}
\PassOptionsToPackage{hyphens}{url}
\PassOptionsToPackage{dvipsnames,svgnames,x11names}{xcolor}
\documentclass[
  10pt,
]{article}
\usepackage{xcolor}
\usepackage[margin=1cm,top=1cm,bottom=1cm,left=1cm,right=1cm,includeheadfoot]{geometry}
\usepackage{amsmath,amssymb}
\setcounter{secnumdepth}{3}
\usepackage{iftex}
\ifPDFTeX
  \usepackage[T1]{fontenc}
  \usepackage[utf8]{inputenc}
  \usepackage{textcomp} % provide euro and other symbols
\else % if luatex or xetex
  \usepackage{unicode-math} % this also loads fontspec
  \defaultfontfeatures{Scale=MatchLowercase}
  \defaultfontfeatures[\rmfamily]{Ligatures=TeX,Scale=1}
\fi
\usepackage{lmodern}
\ifPDFTeX\else
  % xetex/luatex font selection
  \setmainfont[]{DejaVu Serif}
  \setmonofont[]{DejaVu Sans Mono}
\fi
% Use upquote if available, for straight quotes in verbatim environments
\IfFileExists{upquote.sty}{\usepackage{upquote}}{}
\IfFileExists{microtype.sty}{% use microtype if available
  \usepackage[]{microtype}
  \UseMicrotypeSet[protrusion]{basicmath} % disable protrusion for tt fonts
}{}
\usepackage{setspace}
\makeatletter
\@ifundefined{KOMAClassName}{% if non-KOMA class
  \IfFileExists{parskip.sty}{%
    \usepackage{parskip}
  }{% else
    \setlength{\parindent}{0pt}
    \setlength{\parskip}{6pt plus 2pt minus 1pt}}
}{% if KOMA class
  \KOMAoptions{parskip=half}}
\makeatother
\usepackage{listings}
\newcommand{\passthrough}[1]{#1}
\lstset{defaultdialect=[5.3]Lua}
\lstset{defaultdialect=[x86masm]Assembler}
\usepackage{graphicx}
\makeatletter
\newsavebox\pandoc@box
\newcommand*\pandocbounded[1]{% scales image to fit in text height/width
  \sbox\pandoc@box{#1}%
  \Gscale@div\@tempa{\textheight}{\dimexpr\ht\pandoc@box+\dp\pandoc@box\relax}%
  \Gscale@div\@tempb{\linewidth}{\wd\pandoc@box}%
  \ifdim\@tempb\p@<\@tempa\p@\let\@tempa\@tempb\fi% select the smaller of both
  \ifdim\@tempa\p@<\p@\scalebox{\@tempa}{\usebox\pandoc@box}%
  \else\usebox{\pandoc@box}%
  \fi%
}
% Set default figure placement to htbp
\def\fps@figure{htbp}
\makeatother
\setlength{\emergencystretch}{3em} % prevent overfull lines
\providecommand{\tightlist}{%
  \setlength{\itemsep}{0pt}\setlength{\parskip}{0pt}}
% Enable graphics inclusion and ensure figure numbering works
\usepackage{graphicx}
\renewcommand{\figurename}{Figure}

% Configure fonts for Unicode support with fallbacks
\usepackage{newunicodechar}
\newunicodechar{⁴}{\textsuperscript{4}}
\newunicodechar{₄}{\textsubscript{4}}

% Math notation macros
\newcommand{\norm}[1]{\lVert #1\rVert}
\newcommand{\abs}[1]{\lvert #1\rvert}

% Enhanced code block styling for better contrast and readability
\usepackage{fancyvrb}
\usepackage{xcolor}
\usepackage{listings}

% Define custom colors for code blocks
\definecolor{codebg}{RGB}{245, 245, 245}      % Light gray background
\definecolor{codeborder}{RGB}{200, 200, 200}  % Medium gray border
\definecolor{codefg}{RGB}{50, 50, 50}         % Dark gray text

% Configure Verbatim environment for inline code
\DefineVerbatimEnvironment{Verbatim}{Verbatim}{%
    fontsize=\small,
    frame=single,
    framerule=0.5pt,
    framesep=3pt,
    rulecolor=\color{codeborder},
    bgcolor=\color{codebg},
    fgcolor=\color{codefg}
}

% Configure code block styling
\DefineVerbatimEnvironment{Highlighting}{Verbatim}{%
    fontsize=\footnotesize,
    frame=single,
    framerule=0.5pt,
    framesep=5pt,
    rulecolor=\color{codeborder},
    bgcolor=\color{codebg},
    fgcolor=\color{codefg}
}

% Style inline code with \texttt
\renewcommand{\texttt}[1]{%
    \colorbox{codebg}{\color{codefg}\ttfamily #1}%
}

% Configure listings package for code blocks
\lstset{
    backgroundcolor=\color{codebg},
    basicstyle=\footnotesize\ttfamily\color{codefg},
    breakatwhitespace=false,
    breaklines=true,
    captionpos=b,
    commentstyle=\color{codefg},
    deletekeywords={...},
    escapeinside={\%*}{*)},
    extendedchars=true,
    frame=single,
    framerule=0.5pt,
    framesep=5pt,
    keepspaces=true,
    keywordstyle=\color{codefg},
    language=Python,
    morekeywords={*,...},
    numbers=left,
    numbersep=5pt,
    numberstyle=\tiny\color{codefg},
    rulecolor=\color{codeborder},
    showspaces=false,
    showstringspaces=false,
    showtabs=false,
    stepnumber=1,
    stringstyle=\color{codefg},
    tabsize=2,
    title=\lstname
}

% Override any Pandoc default lstset configurations
\AtBeginDocument{
    \lstset{
        backgroundcolor=\color{codebg},
        basicstyle=\footnotesize\ttfamily\color{codefg},
        frame=single,
        framerule=0.5pt,
        framesep=5pt,
        rulecolor=\color{codeborder},
        numbers=left,
        numbersep=5pt,
        numberstyle=\tiny\color{codefg}
    }
}

% Configure hyperref colors consistently
\AtBeginDocument{
% Override pandoc's hidelinks setting with consistent options
\hypersetup{
    colorlinks=true,
    allcolors=red,
    linkcolor=red,
    urlcolor=red,
    citecolor=red,
    filecolor=red,
    menucolor=red,
    linktoc=all
}
}

% Simple page break support for document structure
\usepackage{bookmark}
\IfFileExists{xurl.sty}{\usepackage{xurl}}{} % add URL line breaks if available
\urlstyle{same}
% fallback for those not using the hyperref driver hyperxmp:
\makeatletter
\@ifundefined{xmpquote}{\newcommand{\xmpquote}[1]{#1}}{}
\makeatother
\hypersetup{
  colorlinks=true,
  linkcolor={red},
  filecolor={red},
  citecolor={red},
  urlcolor={red},
  pdfcreator={LaTeX via pandoc}}

\title{Static IVM Field Learning}
\author{Daniel Ari Friedman\\ ORCID: 0000-0001-6232-9096\\ Email: daniel@activeinference.institute\\ DOI: 10.5281/zenodo.16887791}
\date{October 07, 2026}

\begin{document}
\maketitle

{
\hypersetup{linkcolor=black}
\setcounter{tocdepth}{3}
\tableofcontents
}
\setstretch{1.0}
\section{Static IVM Field Learning}\label{static-ivm-field-learning}

\subsection{Overview}\label{overview}

This section develops machine learning on a static synergetic geometry:
scalar fields defined over the Isotropic Vector Matrix (IVM) lattice,
learned from sparse noisy observations by Laplacian-regularized
kernel-weighted least squares on the lattice graph. All methods live in
the \passthrough{\lstinline!ivm\_field.py!} module
(\passthrough{\lstinline!quadray.IVMField!},
\passthrough{\lstinline!quadray\_shell\_norm!},
\passthrough{\lstinline!shell\_sites!},
\passthrough{\lstinline!fit\_geometry!}) and reuse the quadray machinery
of \passthrough{\lstinline!quadray.py!} described in
\href{03_quadray_methods.md}{Quadray Methods}. The learner uses
\passthrough{\lstinline!numpy!} only --- no machine-learning frameworks
--- and is fully deterministic for a fixed observation set.

\subsection{The IVM Lattice in Quadray
Coordinates}\label{the-ivm-lattice-in-quadray-coordinates}

Quadray coordinates represent IVM close-packed sphere centers as
non-negative integer quadrays normalized so the minimum component is
zero (see \href{03_quadray_methods.md}{Quadray Methods}). Two facts
organize the lattice. First, normalization adds or subtracts
\((k,k,k,k)\), so the component sum is a projective invariant modulo 4.
Second, the IVM sites are exactly one of the four residue classes:

\begin{equation}
\label{eq:ivm-site-condition}
q \text{ is an IVM site} \quad\Longleftrightarrow\quad \textstyle\sum_i q_i \equiv 0 \pmod{4},
\end{equation}

implemented as \passthrough{\lstinline!is\_ivm\_site()!} in
\passthrough{\lstinline!ivm\_field.py!}. The remaining three cosets are
the octahedral (sum \(\equiv 2\)) and tetrahedral (sum \(\equiv 1, 3\))
voids of the packing --- lattice points of the ambient grid, but not
sphere centers.

The shell norm of a site is the \(L^1\) magnitude of its sum-zero
(Coxeter.4D hyperplane) representative. With \(s = \sum_i q_i\):

\begin{equation}
\label{eq:ivm-shell-norm}
N(q) \;=\; \sum_{i=1}^{4} \left|\, q_i - \tfrac{s}{4} \,\right| \;=\; 2k, \qquad k \in \mathbb{Z}_{\geq 0},
\end{equation}

computed by \passthrough{\lstinline!quadray\_shell\_norm()!}. Shell
\(k\) of the lattice is the set of sites with \(N(q) = 2k\);
\passthrough{\lstinline!shell\_sites(k)!} enumerates it by scanning the
bounding box \([0, 2k]^4\), filtering on the membership condition
\eqref{eq:ivm-site-condition} and the norm \eqref{eq:ivm-shell-norm},
and sorting lexicographically --- a deterministic site order. The shell
cardinalities are the cuboctahedral numbers:

\begin{equation}
\label{eq:ivm-shell-cardinality}
\bigl|\, \{ q : N(q) = 2k \} \,\bigr| \;=\; 10k^2 + 2, \qquad k \geq 1,
\end{equation}

giving the sequence 1, 12, 42, 92, 162, \ldots{} for
\(k = 0, 1, 2, 3, 4\) --- the center plus the cuboctahedral numbers,
with shell 1 the twelve-around-one vector equilibrium of
\href{03_quadray_methods.md}{Quadray Methods}.
\passthrough{\lstinline!shell\_cardinalities()!} counts shells of an
enumerated ball, and the test suite pins the sequence exactly.

\subsection{The Field Model}\label{the-field-model}

\passthrough{\lstinline!IVMField.lattice\_ball(radius)!} stores a scalar
field over the lattice ball \(B_R = \{ q : N(q) \leq 2R \}\) as a
one-dimensional \passthrough{\lstinline!numpy!} array keyed by the
deterministic site index of
\passthrough{\lstinline!shell\_ball\_sites()!} (shell-major,
lexicographic within shell), with a dictionary mapping each normalized
site to its array position. Two lattice-graph ingredients drive
learning. The adjacency structure uses the twelve IVM neighbor moves ---
the permutations of \((2,1,1,0)\) collected in
\passthrough{\lstinline!IVM\_NEIGHBOR\_STEPS!} --- so two sites in the
ball are graph-adjacent exactly when one step of close packing separates
them. The graph Laplacian \(L\) (method
\passthrough{\lstinline!IVMField.\_laplacian()!}) is the symmetric
difference operator on that adjacency:

\begin{equation}
\label{eq:ivm-laplacian}
(L f)_i \;=\; \deg(i)\, f_i \;-\; \sum_{j \sim i} f_j .
\end{equation}

\subsection{Learning: Laplacian-Regularized Kernel-Weighted Least
Squares}\label{learning-laplacian-regularized-kernel-weighted-least-squares}

Observations \(y_j\) arrive at a sampled subset \(\Omega\) of sites.
Graph distances \(d(\cdot,\cdot)\) are hop counts from a multi-source
BFS (\passthrough{\lstinline!\_multi\_source\_distances()!} in
\passthrough{\lstinline!ivm\_field.py!}), and the Gaussian kernel over
hops with width \(\tau\) (\passthrough{\lstinline!kernel\_width!}) is:

\begin{equation}
\label{eq:ivm-kernel}
K(d) \;=\; \exp\!\bigl( -(d/\tau)^2 \bigr), \qquad d(i,j) = \text{hop distance}.
\end{equation}

\passthrough{\lstinline!IVMField.learn()!} minimizes a two-regime
objective over the ball, with Laplacian regularization weight
\(\lambda \geq 0\):

\begin{equation}
\label{eq:ivm-objective}
\min_{f} \;\; \sum_{i \in \Omega} \bigl( f_i - y_i \bigr)^2
\;+\; \sum_{i \notin \Omega} c_i \bigl( f_i - t_i \bigr)^2
\;+\; \lambda \sum_{i \sim j} \bigl( f_i - f_j \bigr)^2 ,
\end{equation}

with three data-fidelity terms: observed sites are pinned to their data
with unit weight (no self-smoothing of real observations); unobserved
sites carry a kernel confidence weight and a Nadaraya--Watson kernel
target,

\begin{equation}
\label{eq:ivm-confidence}
c_i \;=\; \max\bigl( K\bigl(\min_{j \in \Omega} d(i,j)\bigr),\ \varepsilon \bigr),
\qquad
t_i \;=\; \frac{\sum_{j \in \Omega} K(d(i,j))\, y_j}{\sum_{j \in \Omega} K(d(i,j))},
\end{equation}

where \(\varepsilon\) --- the \passthrough{\lstinline!\_WEIGHT\_FLOOR!}
constant, \(10^{-12}\), in \passthrough{\lstinline!ivm\_field.py!} ---
is a numerical floor keeping the normal-equation matrix of
\eqref{eq:ivm-objective} positive definite on the connected lattice
ball. The normal equations of \eqref{eq:ivm-objective} are

\begin{equation}
\label{eq:ivm-normal-equations}
\bigl( C + \lambda L \bigr) f \;=\; C\, t ,
\end{equation}

with \(C = \mathrm{diag}(c_i)\) and \(t\) the target vector;
\passthrough{\lstinline!IVMField.learn()!} solves them with a dense
\passthrough{\lstinline!numpy.linalg.solve!}. The Laplacian term
propagates field structure from observed sites across the lattice graph,
and the confidence weighting degrades gracefully to pure Laplacian
propagation as observation density falls. Because observed rows carry
unit weight, the estimator interpolates exact data as \(\lambda \to 0\)
and denoises noisy data for moderate \(\lambda\) --- both behaviors are
verified numerically in
\passthrough{\lstinline!tests/test\_ivm\_field.py!} with fixed seeds.

Two structural facts justify the estimator. A linear field in the
embedded XYZ coordinates (the \passthrough{\lstinline!to\_xyz()!} image
of \passthrough{\lstinline!quadray.py!}) is harmonic on the IVM graph
--- the twelve neighbor moves sum to zero --- so the harmonic extension
implied by \eqref{eq:ivm-normal-equations} reproduces it exactly from
boundary data (asserted to \(10^{-5}\) in the test suite). And for noisy
smooth fields, the learned field tracks the kernel-weighted local mean,
which averages observation noise across graph neighborhoods:

\begin{equation}
\label{eq:ivm-mse}
\mathrm{MSE}(f) \;=\; \frac{1}{|B_R|} \sum_{i \in B_R} \bigl( f_i - f_i^{\text{truth}} \bigr)^2 ,
\end{equation}

the quantity returned by \passthrough{\lstinline!IVMField.score()!}.

\subsection{Recovering Tetrahedral Geometry from Noisy
Points}\label{recovering-tetrahedral-geometry-from-noisy-points}

\passthrough{\lstinline!fit\_geometry()!} recovers the orientation and
scale of a tetrahedron from noisy 3D point clouds, using the quadray
basis of \passthrough{\lstinline!quadray.py!}. The four canonical vertex
images are \(v_i = \texttt{to\_xyz}(e_i)\), where \(e_i\) are the unit
quadray axes mapped through the 3\$\times\$4 embedding matrix of
\passthrough{\lstinline!quadray.py!}. Given labeled observations
\(p^{(i)}_m \approx R\, v_i + o\) (noisy samples of vertex \(i\) under
rotation \(R\) and offset \(o\)), vertex centroids \(c_i\) are formed
and the linear map \(G\) is fit in closed form:

\begin{equation}
\label{eq:ivm-fit}
G \;=\; \arg\min_{M \in \mathbb{R}^{3\times 3}} \sum_{i=1}^{4}
\bigl\| M v_i - c_i \bigr\|^2 ,
\qquad
\text{scale} \;=\; \operatorname{sign}\bigl(\det G\bigr)\, \bigl|\det G\bigr|^{1/3}.
\end{equation}

The least-squares solve uses
\passthrough{\lstinline!numpy.linalg.lstsq!} on the transpose system, so
noisy multi-sample vertices average out; the residual reported by
\passthrough{\lstinline!TetrahedronFit.residual!} is the RMS
vertex-centroid mismatch. The test suite recovers a rotated, scaled
tetrahedron to \(10^{-6}\) from \(\sigma =
10^{-7}\) noise, verifies the signed scale for reflected (negative
determinant) fits, and checks that doubling the embedding halves the
recovered matrix --- orientation and scale separate cleanly.

\subsection{Demonstration}\label{demonstration}

The script
\passthrough{\lstinline!quadmath/scripts/ivm\_field\_demo.py!} (run with
\passthrough{\lstinline!MPLBACKEND=Agg!}, seed 12) renders the pipeline
on the radius-3 ball (147 lattice sites): the synthetic field (a
harmonic linear part plus a weak quadratic bowl), the noisy observations
at 74 of the 147 sites, and the learned field. The learned
reconstruction beats the raw observation noise (reconstruction MSE
0.0558 against observation-noise MSE 0.0720, printed by the demo):

\begin{figure}
\centering
\pandocbounded{\includegraphics[keepaspectratio,alt={Static IVM field learning on the radius-3 IVM lattice ball (147 sites, shown at their XYZ embedding from the quadray coordinates; axes in embedding units). Panel A: the synthetic ground-truth field --- a harmonic linear part plus a weak quadratic bowl --- with dot color giving field value on the shared color bar (right). Panel B: the noisy observations at 74 of the 147 sites (seed 12). Panel C: the learned field from Laplacian-regularized kernel-weighted least squares (Eqs. --); reconstruction MSE 0.0558 against observation-noise MSE 0.0720. Reproduce with quadmath/scripts/ivm\_field\_demo.py (MPLBACKEND=Agg, seed 12).}]{../output/figures/ivm_field_demo.png}}
\caption{Static IVM field learning on the radius-3 IVM lattice ball (147
sites, shown at their XYZ embedding from the quadray coordinates; axes
in embedding units). Panel A: the synthetic ground-truth field --- a
harmonic linear part plus a weak quadratic bowl --- with dot color
giving field value on the shared color bar (right). Panel B: the noisy
observations at 74 of the 147 sites (seed 12). Panel C: the learned
field from Laplacian-regularized kernel-weighted least squares (Eqs.
\eqref{eq:ivm-objective}--\eqref{eq:ivm-normal-equations});
reconstruction MSE 0.0558 against observation-noise MSE 0.0720.
Reproduce with
\protect\passthrough{\lstinline!quadmath/scripts/ivm\_field\_demo.py!}
(\protect\passthrough{\lstinline!MPLBACKEND=Agg!}, seed 12).}
\end{figure}

\subsection{Cross-References}\label{cross-references}

\begin{itemize}
\tightlist
\item
  Coordinate foundations and the twelve-around-one shell:
  \href{03_quadray_methods.md}{Quadray Methods}
\item
  Optimization methods on tetrahedral lattices:
  \href{04_optimization_in_4d.md}{Optimization in 4D}
\item
  Applications and generalizations: \href{05_extensions.md}{Extensions}
\end{itemize}

\end{document}
